\documentclass[12pt,a4paper]{article}
\usepackage{../../../myconfig}

\begin{document}

\titlebox{Chủ đề: Oxyz}{Phương trình mặt phẳng}{Oxyz: Phương trình mặt phẳng}

% =========================================================================
% DẠNG 1: XÁC ĐỊNH VÉCTƠ PHÁP TUYẾN CỦA MẶT PHẲNG (CÂU 1 -> 3)
% =========================================================================
\questionLinesManual{0}{1}{%
Trong không gian $Oxyz$, cho mặt phẳng $(\alpha)\colon 3x + 2y - 4z + 1 = 0$. Véctơ nào dưới đây là một véctơ pháp tuyến của $(\alpha)$?
}{%
\begin{tasks}(4)
    \task $\vec{n}_1 = (3; 2; 4)$
    \task $\vec{n}_2 = (3; 2; -4)$
    \task $\vec{n}_3 = (2; -4; 1)$
    \task $\vec{n}_4 = (3; -4; 1)$
\end{tasks}
}

\questionLinesManual{0}{2}{%
Trong không gian $Oxyz$, cho mặt phẳng $(P)\colon 2x + 4y - z + 3 = 0$. Véctơ nào dưới đây là một véctơ pháp tuyến của $(P)$?
}{%
\begin{tasks}(4)
    \task $\vec{n}_1 = (2; 4; -1)$
    \task $\vec{n}_2 = (2; -4; 1)$
    \task $\vec{n}_3 = (-2; 4; 1)$
    \task $\vec{n}_4 = (2; 4; 1)$
\end{tasks}
}

\questionLinesManual{0}{3}{%
Trong không gian $Oxyz$, cho mặt phẳng $(P)\colon x + y - \dfrac{z}{2} = 1$. Một véctơ pháp tuyến của mặt phẳng $(P)$ là:
}{%
\begin{tasks}(4)
    \task $\vec{n} = (1; 1; 2)$
    \task $\vec{n} = (2; 2; -1)$
    \task $\vec{n} = (1; 1; -2)$
    \task $\vec{n} = (2; 2; 1)$
\end{tasks}
}

% =========================================================================
% DẠNG 2: BẪY VTPT TRONG PHƯƠNG TRÌNH ĐOẠN CHẮN (CÂU 4 -> 6)
% =========================================================================
\questionLinesManual{0}{4}{%
Trong không gian $Oxyz$, cho mặt phẳng $(P)\colon \dfrac{x}{2} + \dfrac{y}{-3} + \dfrac{z}{4} = 1$. Một véctơ pháp tuyến của mặt phẳng $(P)$ có toạ độ là:
}{%
\begin{tasks}(4)
    \task $\vec{n}_1 = (2; -3; 4)$
    \task $\vec{n}_2 = (6; -4; 3)$
    \task $\vec{n}_3 = (6; 4; 3)$
    \task $\vec{n}_4 = (3; -2; 6)$
\end{tasks}
}

\questionLinesManual{0}{5}{%
Trong không gian $Oxyz$, một véctơ pháp tuyến của mặt phẳng $(P)\colon \dfrac{x}{-1} + \dfrac{y}{2} + \dfrac{z}{-3} = 1$ có toạ độ là:
}{%
\begin{tasks}(4)
    \task $\vec{n}_1 = (-1; 2; -3)$
    \task $\vec{n}_2 = (6; -3; 2)$
    \task $\vec{n}_3 = (6; 3; -2)$
    \task $\vec{n}_4 = (-6; -3; 2)$
\end{tasks}
}

\questionLinesManual{0}{6}{%
Trong không gian $Oxyz$, cho mặt phẳng $(\alpha)\colon \dfrac{x}{3} - \dfrac{y}{2} + z = 1$. Véctơ nào dưới đây là một véctơ pháp tuyến của mặt phẳng $(\alpha)$?
}{%
\begin{tasks}(4)
    \task $\vec{n}_1 = (3; -2; 1)$
    \task $\vec{n}_2 = (2; -3; 1)$
    \task $\vec{n}_3 = (2; -3; 6)$
    \task $\vec{n}_4 = (2; 3; 6)$
\end{tasks}
}

% =========================================================================
% DẠNG 3: BẪY VTPT KHUYẾT BIẾN / RÚT GỌN TỌA ĐỘ (CÂU 7 -> 9)
% =========================================================================
\questionLinesManual{0}{7}{%
Trong không gian $Oxyz$, cho mặt phẳng $(P)\colon 3x - z + 2 = 0$. Véctơ nào dưới đây là một véctơ pháp tuyến của $(P)$?
}{%
\begin{tasks}(4)
    \task $\vec{n} = (3; 0; 2)$
    \task $\vec{n} = (3; -1; 0)$
    \task $\vec{n} = (3; 0; -1)$
    \task $\vec{n} = (-1; 0; -1)$
\end{tasks}
}

\questionLinesManual{0}{8}{%
Trong không gian $Oxyz$, cho mặt phẳng $(P)\colon 3x - 2z + 5 = 0$. Véctơ nào dưới đây là một véctơ pháp tuyến của $(P)$?
}{%
\begin{tasks}(4)
    \task $\vec{n}_1 = (3; -2; 0)$
    \task $\vec{n}_2 = (3; 0; -2)$
    \task $\vec{n}_3 = (3; -2; 5)$
    \task $\vec{n}_4 = (3; 0; 2)$
\end{tasks}
}

\questionLinesManual{0}{9}{%
Trong không gian $Oxyz$, cho mặt phẳng $(P)\colon 2x - 4y + 8z - 1 = 0$. Véctơ nào sau đây là một véctơ pháp tuyến của $(P)$?
}{%
\begin{tasks}(4)
    \task $\vec{n} = (1; 2; 4)$
    \task $\vec{n} = (1; -2; 4)$
    \task $\vec{n} = (2; -4; 4)$
    \task $\vec{n} = (2; -4; 1)$
\end{tasks}
}

% =========================================================================
% DẠNG 4: VTPT VÀ PHƯƠNG TRÌNH CỦA MẶT PHẲNG TỌA ĐỘ (CÂU 10 -> 12)
% =========================================================================
\questionLinesManual{0}{10}{%
Trong không gian $Oxyz$, véctơ nào dưới đây là một véctơ pháp tuyến của mặt phẳng tọa độ $(Oxy)$?
}{%
\begin{tasks}(4)
    \task $\vec{i} = (1; 0; 0)$
    \task $\vec{j} = (0; 1; 0)$
    \task $\vec{k} = (0; 0; 1)$
    \task $\vec{m} = (1; 1; 1)$
\end{tasks}
}

\questionLinesManual{0}{11}{%
Trong không gian $Oxyz$, mặt phẳng $(Oxz)$ có một véctơ pháp tuyến là:
}{%
\begin{tasks}(4)
    \task $\vec{n} = (0; 1; 0)$
    \task $\vec{n} = (1; 0; 0)$
    \task $\vec{n} = (0; 0; 1)$
    \task $\vec{n} = (1; 0; 1)$
\end{tasks}
}

\questionLinesManual{0}{12}{%
Trong không gian $Oxyz$, phương trình nào sau đây là phương trình của mặt phẳng $(Oxz)$?
}{%
\begin{tasks}(4)
    \task $x = 0$
    \task $y = 0$
    \task $z = 0$
    \task $x + z = 0$
\end{tasks}
}

% =========================================================================
% DẠNG 5: CẶP VTCP VÀ TÍCH CÓ HƯỚNG TÌM VTPT (CÂU 13 -> 15)
% =========================================================================
\questionLinesManual{0}{13}{%
Trong không gian $Oxyz$, mặt phẳng $(P)$ nhận hai véctơ $\vec{u}_1 = (1; 2; -1)$ và $\vec{u}_2 = (0; -1; 3)$ làm cặp véctơ chỉ phương. Một véctơ pháp tuyến của $(P)$ là:
}{%
\begin{tasks}(4)
    \task $\vec{n} = (5; -3; -1)$
    \task $\vec{n} = (5; 3; -1)$
    \task $\vec{n} = (-5; 3; 1)$
    \task $\vec{n} = (5; -3; 1)$
\end{tasks}
}

\questionLinesManual{0}{14}{%
Trong không gian $Oxyz$, cho ba điểm $A(1; 1; 1)$, $B(1; 0; 1)$, $C(0; 0; 1)$. Biết mặt phẳng $(ABC)$ có một véctơ pháp tuyến dạng $\vec{n} = (a; b; -1)$ với $a, b \in \mathbb{R}$. Giá trị của $a + b$ bằng:
}{%
\begin{tasks}(4)
    \task $2$
    \task $1$
    \task $0$
    \task $-1$
\end{tasks}
}

\questionLinesManual{0}{15}{%
Trong không gian $Oxyz$, cho ba điểm $A(2; -1; 3)$, $B(4; 0; 1)$ và $C(-10; 5; 3)$. Véctơ nào dưới đây là một véctơ pháp tuyến của mặt phẳng $(ABC)$?
}{%
\begin{tasks}(4)
    \task $\vec{n}_1 = (1; 2; 0)$
    \task $\vec{n}_2 = (1; 2; 2)$
    \task $\vec{n}_3 = (1; 8; 2)$
    \task $\vec{n}_4 = (1; -2; 2)$
\end{tasks}
}

% =========================================================================
% DẠNG 6: KIỂM TRA ĐIỂM THUỘC / KHÔNG THUỘC MẶT PHẲNG (CÂU 16 -> 18)
% =========================================================================
\questionLinesManual{0}{16}{%
Trong không gian $Oxyz$, cho mặt phẳng $(\alpha)\colon x + y + z - 6 = 0$. Điểm nào dưới đây không thuộc $(\alpha)$?
}{%
\begin{tasks}(4)
    \task $Q(3; 3; 0)$
    \task $N(2; 2; 2)$
    \task $P(1; 2; 3)$
    \task $M(1; -1; 1)$
\end{tasks}
}

\questionLinesManual{0}{17}{%
Trong không gian $Oxyz$, cho mặt phẳng $(P)\colon x - 2y + z - 5 = 0$. Điểm nào dưới đây thuộc $(P)$?
}{%
\begin{tasks}(4)
    \task $P(0; 0; -5)$
    \task $M(1; 1; 6)$
    \task $Q(2; -1; 5)$
    \task $N(-5; 0; 0)$
\end{tasks}
}

\questionLinesManual{0}{18}{%
Trong không gian $Oxyz$, cho mặt phẳng $(P)\colon 2x - y + z - 5 = 0$. Điểm nào sau đây thuộc $(P)$?
}{%
\begin{tasks}(4)
    \task $M(1; -1; 0)$
    \task $N(1; -1; 2)$
    \task $P(1; -1; 4)$
    \task $Q(1; -1; 3)$
\end{tasks}
}

% =========================================================================
% DẠNG 7: BẪY MẶT TỌA ĐỘ VÀ ĐỐI XỨNG / HÌNH CHIẾU CƠ BẢN (CÂU 19 -> 21)
% =========================================================================
\questionLinesManual{0}{19}{%
Trong không gian $Oxyz$, hình chiếu vuông góc của điểm $A(2; 3; -2)$ trên mặt phẳng tọa độ $(Oyz)$ có toạ độ là:
}{%
\begin{tasks}(4)
    \task $(-2; 3; -2)$
    \task $(2; 3; 0)$
    \task $(2; 0; 0)$
    \task $(0; 3; -2)$
\end{tasks}
}

\questionLinesManual{0}{20}{%
Trong không gian $Oxyz$, tìm tọa độ điểm $M'$ là điểm đối xứng của điểm $A(1; -2; 5)$ qua mặt phẳng tọa độ $(Oxy)$.
}{%
\begin{tasks}(4)
    \task $(1; -2; -5)$
    \task $(-1; 2; 5)$
    \task $(1; 2; -5)$
    \task $(1; 2; 5)$
\end{tasks}
}

\questionLinesManual{0}{21}{%
Trong không gian $Oxyz$, điểm nào dưới đây thuộc mặt phẳng tọa độ $(Oyz)$?
}{%
\begin{tasks}(4)
    \task $M(2; 0; 0)$
    \task $N(0; -3; 1)$
    \task $P(1; 0; 2)$
    \task $Q(1; 1; 0)$
\end{tasks}
}

% =========================================================================
% DẠNG 8: VIẾT MP QUA 1 ĐIỂM VÀ CÓ VTPT CHO TRƯỚC (CÂU 22 -> 24)
% =========================================================================
\questionLinesManual{0}{22}{%
Trong không gian $Oxyz$, cho mặt phẳng $(P)$ đi qua điểm $A(1; 0; -2)$ và có một véctơ pháp tuyến $\vec{n} = (1; -1; 2)$. Phương trình mặt phẳng $(P)$ là:
}{%
\begin{tasks}(2)
    \task $x - y - 2z + 3 = 0$
    \task $x - y + 2z - 3 = 0$
    \task $x - y + 2z + 3 = 0$
    \task $x + y + 2z + 3 = 0$
\end{tasks}
}

\questionLinesManual{0}{23}{%
Trong không gian $Oxyz$, phương trình mặt phẳng đi qua điểm $M(1; -2; 3)$ và có véctơ pháp tuyến $\vec{n} = (2; 0; 1)$ là:
}{%
\begin{tasks}(2)
    \task $2x + z - 5 = 0$
    \task $2x + y = 0$
    \task $x - 2y + 3z - 5 = 0$
    \task $2y + z + 1 = 0$
\end{tasks}
}

\questionLinesManual{0}{24}{%
Trong không gian $Oxyz$, mặt phẳng đi qua điểm $A(2; 1; 0)$ và nhận véctơ $\vec{n} = (3; -1; -1)$ làm véctơ pháp tuyến có phương trình là:
}{%
\begin{tasks}(2)
    \task $3x - y - z + 5 = 0$
    \task $3x - y - z - 5 = 0$
    \task $2x + y - 5 = 0$
    \task $x + 3y + z - 5 = 0$
\end{tasks}
}

% =========================================================================
% DẠNG 9: BẪY VIẾT MP QUA ĐIỂM SONG SONG/VUÔNG GÓC MẶT TỌA ĐỘ (CÂU 25 -> 27)
% =========================================================================
\questionLinesManual{0}{25}{%
Trong không gian $Oxyz$, mặt phẳng đi qua điểm $M(1; -2; 3)$ đồng thời song song với mặt phẳng tọa độ $(Oxy)$ có phương trình là:
}{%
\begin{tasks}(4)
    \task $x - 1 = 0$
    \task $y + 2 = 0$
    \task $z - 3 = 0$
    \task $z + 3 = 0$
\end{tasks}
}

\questionLinesManual{0}{26}{%
Trong không gian $Oxyz$, mặt phẳng $(P)$ đi qua điểm $A(2; 0; -3)$ và vuông góc với trục hoành $Ox$ có phương trình là:
}{%
\begin{tasks}(4)
    \task $x - 2 = 0$
    \task $x + 2 = 0$
    \task $2x - 3z = 0$
    \task $z + 3 = 0$
\end{tasks}
}

\questionLinesManual{0}{27}{%
Trong không gian $Oxyz$, mặt phẳng $(P)$ đi qua điểm $B(1; -4; 2)$ và vuông góc với trục tung $Oy$ có phương trình là:
}{%
\begin{tasks}(4)
    \task $x - 1 = 0$
    \task $y - 4 = 0$
    \task $y + 4 = 0$
    \task $z - 2 = 0$
\end{tasks}
}

% =========================================================================
% DẠNG 10: VIẾT MP QUA 1 ĐIỂM VÀ SONG SONG VỚI MẶT PHẲNG CHO TRƯỚC (CÂU 28 -> 30)
% =========================================================================
\questionLinesManual{0}{28}{%
Trong không gian $Oxyz$, cho điểm $M(2; -1; 4)$ và mặt phẳng $(P)\colon 3x - 2y + z + 1 = 0$. Phương trình của mặt phẳng đi qua $M$ và song song với mặt phẳng $(P)$ là:
}{%
\begin{tasks}(2)
    \task $2x - 2y + 4z - 21 = 0$
    \task $3x - 2y + z + 12 = 0$
    \task $3x - 2y + z - 12 = 0$
    \task $2x - 2y + 4z + 21 = 0$
\end{tasks}
}

\questionLinesManual{0}{29}{%
Trong không gian $Oxyz$, cho mặt phẳng $(P)\colon 2x - y + z - 4 = 0$. Phương trình mặt phẳng $(Q)$ song song với $(P)$ và đi qua điểm $M(1; -1; 3)$ là:
}{%
\begin{tasks}(2)
    \task $2x - y + z + 6 = 0$
    \task $2x + 3y - z + 4 = 0$
    \task $2x - y + z - 6 = 0$
    \task $x - y + 3z - 6 = 0$
\end{tasks}
}

\questionLinesManual{0}{30}{%
Trong không gian $Oxyz$, cho điểm $A(0; -3; 2)$ và mặt phẳng $(P)\colon 2x - y + 3z + 5 = 0$. Mặt phẳng đi qua $A$ và song song với $(P)$ có phương trình là:
}{%
\begin{tasks}(2)
    \task $2x + y + 3z + 3 = 0$
    \task $2x - y + 3z - 9 = 0$
    \task $2x - y + 3z + 9 = 0$
    \task $2x + y + 3z - 3 = 0$
\end{tasks}
}

% =========================================================================
% DẠNG 11: MẶT PHẲNG TRUNG TRỰC CỦA ĐOẠN THẲNG (CÂU 31 -> 33)
% =========================================================================
\questionLinesManual{0}{31}{%
Trong không gian $Oxyz$, cho hai điểm $A(4; 0; 1)$ và $B(-2; 2; 3)$. Mặt phẳng trung trực của đoạn thẳng $AB$ có phương trình là:
}{%
\begin{tasks}(2)
    \task $3x - y - z = 0$
    \task $3x + y + z - 6 = 0$
    \task $x + y + 2z - 6 = 0$
    \task $6x - 2y - 2z - 1 = 0$
\end{tasks}
}

\questionLinesManual{0}{32}{%
Trong không gian $Oxyz$, cho hai điểm $A(1; 3; 0)$ và $B(5; 1; -1)$. Mặt phẳng trung trực của đoạn thẳng $AB$ có phương trình là:
}{%
\begin{tasks}(2)
    \task $x + y + 2z - 3 = 0$
    \task $3x + 2y - z - 14 = 0$
    \task $2x - y - z + 5 = 0$
    \task $2x - y - z - 5 = 0$
\end{tasks}
}

\questionLinesManual{0}{33}{%
Trong không gian $Oxyz$, cho hai điểm $A(1; -1; 2)$ và $B(3; 3; 0)$. Mặt phẳng trung trực của đoạn thẳng $AB$ có phương trình là:
}{%
\begin{tasks}(2)
    \task $x + y - z - 2 = 0$
    \task $x + y - z + 2 = 0$
    \task $x + 2y - z - 3 = 0$
    \task $x + 2y - z + 3 = 0$
\end{tasks}
}

% =========================================================================
% DẠNG 12: VIẾT MP QUA ĐIỂM VUÔNG GÓC VỚI ĐOẠN NỐI (CÂU 34 -> 36)
% =========================================================================
\questionLinesManual{0}{34}{%
Trong không gian $Oxyz$, cho hai điểm $A(0; 1; 1)$ và $B(1; 2; 3)$. Viết phương trình mặt phẳng $(P)$ đi qua $A$ và vuông góc với đoạn thẳng $AB$.
}{%
\begin{tasks}(2)
    \task $x + y + 2z - 3 = 0$
    \task $x + y + 2z - 6 = 0$
    \task $x + 3y + 4z - 7 = 0$
    \task $x + 3y + 4z - 26 = 0$
\end{tasks}
}

\questionLinesManual{0}{35}{%
Trong không gian $Oxyz$, cho ba điểm $A(2; 1; -1)$, $B(-1; 0; 4)$ và $C(0; -2; -1)$. Phương trình mặt phẳng đi qua $A$ và vuông góc với đoạn thẳng $BC$ là:
}{%
\begin{tasks}(2)
    \task $x - 2y - 5z - 5 = 0$
    \task $2x - y + 5z - 5 = 0$
    \task $x - 2y - 5 = 0$
    \task $x - 2y - 5z + 5 = 0$
\end{tasks}
}

\questionLinesManual{0}{36}{%
Trong không gian $Oxyz$, cho hai điểm $A(1; 0; 2)$ và $B(2; -1; 3)$. Viết phương trình mặt phẳng $(P)$ đi qua $A$ và vuông góc với đoạn thẳng $AB$.
}{%
\begin{tasks}(2)
    \task $x - y + z - 3 = 0$
    \task $2x - y + z - 4 = 0$
    \task $-x + 2y + z - 1 = 0$
    \task $x + y + z - 3 = 0$
\end{tasks}
}

% =========================================================================
% DẠNG 13: VIẾT MP QUA 1 ĐIỂM VÀ CÓ CẶP VTCP CHO TRƯỚC (CÂU 37 -> 39)
% =========================================================================
\questionLinesManual{0}{37}{%
Trong không gian $Oxyz$, phương trình mặt phẳng $(P)$ đi qua điểm $A(1; 2; 0)$ và nhận cặp véctơ $\vec{u}_1 = (1; 0; 2)$, $\vec{u}_2 = (-1; 1; 1)$ làm cặp véctơ chỉ phương là:
}{%
\begin{tasks}(2)
    \task $2x + 3y - z - 8 = 0$
    \task $2x - 3y + z + 4 = 0$
    \task $-2x - 3y + z + 8 = 0$
    \task $2x - 3y - z + 4 = 0$
\end{tasks}
}

\questionLinesManual{0}{38}{%
Trong không gian $Oxyz$, viết phương trình mặt phẳng $(\alpha)$ đi qua điểm $M(2; -1; 3)$ đồng thời có cặp véctơ chỉ phương $\vec{a} = (2; 1; -1)$ và $\vec{b} = (0; 3; 2)$.
}{%
\begin{tasks}(2)
    \task $5x - 4y + 6z - 32 = 0$
    \task $5x + 4y - 6z + 12 = 0$
    \task $5x - 4y + 6z + 32 = 0$
    \task $5x - 4y - 6z + 4 = 0$
\end{tasks}
}

\questionLinesManual{0}{39}{%
Trong không gian $Oxyz$, cho điểm $A(1; -1; 5)$ và $B(0; 0; 1)$. Mặt phẳng chứa hai điểm $A, B$ và song song với trục tung $Oy$ có phương trình là:
}{%
\begin{tasks}(4)
    \task $4x - z + 1 = 0$
    \task $2x + z - 3 = 0$
    \task $x - 4z + 2 = 0$
    \task $4x - z - 1 = 0$
\end{tasks}
}

% =========================================================================
% DẠNG 14: VIẾT PHƯƠNG TRÌNH MẶT PHẲNG ĐI QUA 3 ĐIỂM (CÂU 40 -> 42)
% =========================================================================
\questionLinesManual{0}{40}{%
Trong không gian $Oxyz$, cho ba điểm $A(1; -3; 2)$, $B(1; 0; 1)$ và $C(2; 3; 0)$. Viết phương trình mặt phẳng $(ABC)$.
}{%
\begin{tasks}(2)
    \task $3x - y - 3z = 0$
    \task $3x + y + 3z - 6 = 0$
    \task $15x - y - 3z - 12 = 0$
    \task $y + 3z - 3 = 0$
\end{tasks}
}

\questionLinesManual{0}{41}{%
Trong không gian $Oxyz$, cho ba điểm $A(3; -2; -2)$, $B(3; 2; 0)$ và $C(0; 2; 1)$. Phương trình mặt phẳng $(ABC)$ là:
}{%
\begin{tasks}(2)
    \task $2x - 3y + 6z + 12 = 0$
    \task $2x + 3y - 6z - 12 = 0$
    \task $2x - 3y + 6z = 0$
    \task $2x + 3y + 6z + 12 = 0$
\end{tasks}
}

\questionLinesManual{0}{42}{%
Trong không gian $Oxyz$, cho ba điểm $A(2; 3; 5)$, $B(3; 2; 4)$ và $C(4; 1; 2)$. Phương trình mặt phẳng đi qua ba điểm $A, B, C$ là:
}{%
\begin{tasks}(2)
    \task $x + y + 5 = 0$
    \task $x + y - 5 = 0$
    \task $y - z + 2 = 0$
    \task $2x + y - 7 = 0$
\end{tasks}
}

% =========================================================================
% DẠNG 15: MẶT PHẲNG QUA 2 ĐIỂM VUÔNG GÓC VỚI 1 MẶT PHẲNG KHÁC (CÂU 43 -> 45)
% =========================================================================
\questionLinesManual{0}{43}{%
Trong không gian $Oxyz$, cho hai điểm $A(2; 4; 1)$, $B(-1; 1; 3)$ và mặt phẳng $(P)\colon x - 3y + 2z - 5 = 0$. Lập phương trình mặt phẳng $(Q)$ đi qua hai điểm $A, B$ và vuông góc với mặt phẳng $(P)$.
}{%
\begin{tasks}(2)
    \task $2y + 3z - 11 = 0$
    \task $2x - 3y - 11 = 0$
    \task $x - 3y + 2z - 5 = 0$
    \task $3y + 2z - 11 = 0$
\end{tasks}
}

\questionLinesManual{0}{44}{%
Trong không gian $Oxyz$, cho $A(1; 2; -1)$, $B(-1; 0; 1)$ và mặt phẳng $(P)\colon x + 2y - z + 1 = 0$. Viết phương trình mặt phẳng $(Q)$ qua $A, B$ và vuông góc với $(P)$.
}{%
\begin{tasks}(4)
    \task $2x - y + 3 = 0$
    \task $x + z = 0$
    \task $-x + y + z = 0$
    \task $3x - y + z = 0$
\end{tasks}
}

\questionLinesManual{0}{45}{%
Trong không gian $Oxyz$, cho mặt phẳng $(P)$ đi qua hai điểm $A(3; 1; -1)$, $B(2; -1; 4)$ và vuông góc với mặt phẳng $(Q)\colon 2x - y + 3z - 1 = 0$. Phương trình của $(P)$ là:
}{%
\begin{tasks}(2)
    \task $x - 13y - 5z + 5 = 0$
    \task $x - 13y + 5z + 5 = 0$
    \task $x + 13y - 5z + 5 = 0$
    \task $x - 13y - 5z + 12 = 0$
\end{tasks}
}

% =========================================================================
% DẠNG 16: MẶT PHẲNG QUA 1 ĐIỂM VUÔNG GÓC VỚI 2 MẶT PHẲNG (CÂU 46 -> 48)
% =========================================================================
\questionLinesManual{0}{46}{%
Trong không gian $Oxyz$, cho hai mặt phẳng $(\alpha)\colon 3x - 2y + 2z + 7 = 0$ và $(\beta)\colon 5x - 4y + 3z + 1 = 0$. Phương trình mặt phẳng đi qua gốc tọa độ $O$ đồng thời vuông góc với cả $(\alpha)$ và $(\beta)$ là:
}{%
\begin{tasks}(2)
    \task $2x - y - 2z = 0$
    \task $2x - y + 2z = 0$
    \task $2x + y - 2z = 0$
    \task $2x + y - 2z + 1 = 0$
\end{tasks}
}

\questionLinesManual{0}{47}{%
Trong không gian $Oxyz$, phương trình của mặt phẳng $(P)$ đi qua điểm $B(2; 1; -3)$ đồng thời vuông góc với hai mặt phẳng $(Q)\colon x + y + 3z = 0$ và $(R)\colon 2x - y + z = 0$ là:
}{%
\begin{tasks}(2)
    \task $4x + 5y - 3z + 22 = 0$
    \task $4x - 5y - 3z - 12 = 0$
    \task $2x + y - 3z - 14 = 0$
    \task $4x + 5y - 3z - 22 = 0$
\end{tasks}
}

\questionLinesManual{0}{48}{%
Trong không gian $Oxyz$, cho điểm $A(1; 1; 1)$ và hai mặt phẳng $(P)\colon 2x - y + 3z - 1 = 0$, $(Q)\colon y = 0$. Viết phương trình mặt phẳng $(R)$ chứa $A$, vuông góc với cả $(P)$ và $(Q)$.
}{%
\begin{tasks}(2)
    \task $3x - y + 2z - 4 = 0$
    \task $3x + y - 2z - 2 = 0$
    \task $3x - 2z = 0$
    \task $3x - 2z - 1 = 0$
\end{tasks}
}

% =========================================================================
% DẠNG 17: BẪY VIẾT PHƯƠNG TRÌNH ĐOẠN CHẮN (CÂU 49 -> 51)
% =========================================================================
\questionLinesManual{0}{49}{%
Trong không gian $Oxyz$, cho ba điểm $A(3; 0; 0)$, $B(0; 1; 0)$ và $C(0; 0; -2)$. Mặt phẳng $(ABC)$ có phương trình là:
}{%
\begin{tasks}(2)
    \task $\dfrac{x}{3} + \dfrac{y}{-1} + \dfrac{z}{2} = 1$
    \task $\dfrac{x}{3} + \dfrac{y}{1} + \dfrac{z}{-2} = 1$
    \task $\dfrac{x}{3} + \dfrac{y}{1} + \dfrac{z}{2} = 1$
    \task $\dfrac{x}{-3} + \dfrac{y}{1} + \dfrac{z}{2} = 1$
\end{tasks}
}

\questionLinesManual{0}{50}{%
Trong không gian $Oxyz$, cho ba điểm $M(0; 2; 0)$, $N(-3; 0; 0)$, $P(0; 0; 5)$. Phương trình theo đoạn chắn của mặt phẳng $(MNP)$ là:
}{%
\begin{tasks}(2)
    \task $\dfrac{x}{2} + \dfrac{y}{-3} + \dfrac{z}{5} = 1$
    \task $\dfrac{x}{-3} + \dfrac{y}{2} + \dfrac{z}{5} = 0$
    \task $\dfrac{x}{-3} + \dfrac{y}{2} + \dfrac{z}{5} = 1$
    \task $\dfrac{x}{2} - \dfrac{y}{3} + \dfrac{z}{5} = -1$
\end{tasks}
}

\questionLinesManual{0}{51}{%
Trong không gian $Oxyz$, phương trình mặt phẳng đi qua ba điểm $A(0; 0; -2)$, $B(0; 3; 0)$, $C(1; 0; 0)$ là:
}{%
\begin{tasks}(2)
    \task $\dfrac{x}{-2} + \dfrac{y}{3} + \dfrac{z}{1} = 1$
    \task $\dfrac{x}{1} + \dfrac{y}{3} - \dfrac{z}{2} = 1$
    \task $\dfrac{x}{1} + \dfrac{y}{3} + \dfrac{z}{2} = 1$
    \task $\dfrac{x}{1} + \dfrac{y}{3} - \dfrac{z}{2} = 0$
\end{tasks}
}

% =========================================================================
% DẠNG 18: TỌA ĐỘ GIAO ĐIỂM VỚI CÁC TRỤC & MẶT PHẲNG ĐOẠN CHẮN QUA HÌNH CHIẾU (CÂU 52 -> 54)
% =========================================================================
\questionLinesManual{0}{52}{%
Trong không gian $Oxyz$, cho mặt phẳng $(P)\colon 2x - 3y + 6z - 6 = 0$. Mặt phẳng $(P)$ cắt trục tung $Oy$ tại điểm nào dưới đây?
}{%
\begin{tasks}(4)
    \task $B(0; -3; 0)$
    \task $B(0; 2; 0)$
    \task $B(0; -2; 0)$
    \task $B(0; 3; 0)$
\end{tasks}
}

\questionLinesManual{0}{53}{%
Trong không gian $Oxyz$, cho điểm $M(1; 2; 3)$. Gọi $A, B, C$ lần lượt là hình chiếu vuông góc của $M$ lên các trục tọa độ $Ox, Oy, Oz$. Viết phương trình mặt phẳng $(ABC)$.
}{%
\begin{tasks}(2)
    \task $\dfrac{x}{1} + \dfrac{y}{2} + \dfrac{z}{3} = 1$
    \task $\dfrac{x}{1} - \dfrac{y}{2} + \dfrac{z}{3} = 1$
    \task $\dfrac{x}{1} + \dfrac{y}{2} + \dfrac{z}{3} = 0$
    \task $-\dfrac{x}{1} + \dfrac{y}{2} + \dfrac{z}{3} = 1$
\end{tasks}
}

\questionLinesManual{0}{54}{%
Trong không gian $Oxyz$, cho điểm $A(1; 2; -5)$. Gọi $M, N, P$ là hình chiếu của $A$ lên các trục $Ox, Oy, Oz$. Phương trình mặt phẳng $(MNP)$ là:
}{%
\begin{tasks}(2)
    \task $x + \dfrac{y}{2} - \dfrac{z}{5} = 1$
    \task $x + 2y - 5z + 1 = 0$
    \task $x + 2y - 5z = 1$
    \task $x + \dfrac{y}{2} - \dfrac{z}{5} + 1 = 0$
\end{tasks}
}

% =========================================================================
% DẠNG 19: BẪY ĐOẠN CHẮN BẰNG NHAU / TỈ LỆ (CÂU 55 -> 57)
% =========================================================================
\questionLinesManual{0}{55}{%
Trong không gian $Oxyz$, mặt phẳng $(\alpha)$ đi qua điểm $M(5; 4; 3)$ và chắn trên các tia $Ox, Oy, Oz$ các đoạn bằng nhau có phương trình là:
}{%
\begin{tasks}(2)
    \task $x - y + z - 4 = 0$
    \task $x + y + z - 12 = 0$
    \task $5x + 4y + 3z - 50 = 0$
    \task $x - y - z + 2 = 0$
\end{tasks}
}

\questionLinesManual{0}{56}{%
Trong không gian $Oxyz$, mặt phẳng $(P)$ đi qua $M(1; 2; 3)$ và cắt các tia $Ox, Oy, Oz$ lần lượt tại $A, B, C$ sao cho $OA = OB = OC > 0$ có phương trình là:
}{%
\begin{tasks}(2)
    \task $x + y + z + 6 = 0$
    \task $x + y + z - 6 = 0$
    \task $\dfrac{x}{1} + \dfrac{y}{2} + \dfrac{z}{3} = 1$
    \task $x + 2y + 3z - 6 = 0$
\end{tasks}
}

\questionLinesManual{0}{57}{%
Có bao nhiêu mặt phẳng đi qua điểm $M(1; 2; 5)$ và cắt các trục $Ox, Oy, Oz$ lần lượt tại $A, B, C$ sao cho $OA = OB = OC \neq 0$?
}{%
\begin{tasks}(4)
    \task $1$
    \task $2$
    \task $3$
    \task $4$
\end{tasks}
}

% =========================================================================
% DẠNG 20: MẶT PHẲNG ĐOẠN CHẮN LIÊN QUAN TRỰC TÂM / TRỌNG TÂM (CÂU 58 -> 60)
% =========================================================================
\questionLinesManual{0}{58}{%
Trong không gian $Oxyz$, cho điểm $H(1; 2; 3)$. Mặt phẳng $(P)$ đi qua điểm $H$, cắt các tia $Ox, Oy, Oz$ lần lượt tại $A, B, C$ sao cho $H$ là trực tâm của tam giác $ABC$. Phương trình của mặt phẳng $(P)$ là:
}{%
\begin{tasks}(2)
    \task $3x + y + 2z - 11 = 0$
    \task $3x + 2y + z - 10 = 0$
    \task $x + 3y + 2z - 13 = 0$
    \task $x + 2y + 3z - 14 = 0$
\end{tasks}
}

\questionLinesManual{0}{59}{%
Trong không gian $Oxyz$, cho mặt phẳng $(P)$ cắt ba trục $Ox, Oy, Oz$ lần lượt tại $A, B, C$ sao cho tam giác $ABC$ có trọng tâm là $G(-1; -3; 2)$. Phương trình mặt phẳng $(P)$ là:
}{%
\begin{tasks}(2)
    \task $6x + 2y - 3z + 18 = 0$
    \task $\dfrac{x}{3} + \dfrac{y}{9} - \dfrac{z}{6} = 1$
    \task $\dfrac{x}{-3} + \dfrac{y}{-9} + \dfrac{z}{6} = 0$
    \task $\dfrac{x}{-1} + \dfrac{y}{-3} + \dfrac{z}{2} = 1$
\end{tasks}
}

\questionLinesManual{0}{60}{%
Trong không gian $Oxyz$, cho điểm $G(1; 4; 3)$. Mặt phẳng nào sau đây cắt các trục $Ox, Oy, Oz$ lần lượt tại $A, B, C$ sao cho $G$ là trọng tâm tứ diện $OABC$?
}{%
\begin{tasks}(2)
    \task $\dfrac{x}{3} + \dfrac{y}{12} + \dfrac{z}{9} = 1$
    \task $12x + 3y + 4z - 48 = 0$
    \task $\dfrac{x}{4} + \dfrac{y}{16} + \dfrac{z}{12} = 0$
    \task $12x + 3y + 4z = 0$
\end{tasks}
}

% =========================================================================
% DẠNG 21: HAI MẶT PHẲNG SONG SONG / VUÔNG GÓC CƠ BẢN (CÂU 61 -> 63)
% =========================================================================
\questionLinesManual{0}{61}{%
Trong không gian $Oxyz$, cho mặt phẳng $(\alpha)\colon x + y + 2z + 2 = 0$. Mặt phẳng nào sau đây song song với $(\alpha)$?
}{%
\begin{tasks}(2)
    \task $(S)\colon x + y + 2z - 1 = 0$
    \task $(Q)\colon x + y - 2z - 2 = 0$
    \task $(P)\colon x - y + 2z - 2 = 0$
    \task $(R)\colon x + y - 2z + 1 = 0$
\end{tasks}
}

\questionLinesManual{0}{62}{%
Trong không gian $Oxyz$, mặt phẳng $(P)\colon 2x + y + z - 2 = 0$ vuông góc với mặt phẳng nào dưới đây?
}{%
\begin{tasks}(2)
    \task $2x - y - z - 2 = 0$
    \task $x - y - z - 2 = 0$
    \task $x + y + z - 2 = 0$
    \task $2x + y + z - 2 = 0$
\end{tasks}
}

\questionLinesManual{0}{63}{%
Trong không gian $Oxyz$, cho mặt phẳng $(\alpha)\colon 2x - 3y + 4z - 5 = 0$. Mặt phẳng nào dưới đây vuông góc với $(\alpha)$?
}{%
\begin{tasks}(2)
    \task $4y + 3z = 0$
    \task $3x - 2y = 0$
    \task $4y - 3z = 0$
    \task $3x + 4y - 1 = 0$
\end{tasks}
}

% =========================================================================
% DẠNG 22: TÌM THAM SỐ ĐỂ HAI MẶT PHẲNG VUÔNG GÓC / SONG SONG (CÂU 64 -> 66)
% =========================================================================
\questionLinesManual{0}{64}{%
Trong không gian $Oxyz$, cho hai mặt phẳng $(P)\colon x - 2y + 2z - 3 = 0$ và $(Q)\colon mx + y - 2z + 1 = 0$. Với giá trị nào của $m$ thì hai mặt phẳng đó vuông góc với nhau?
}{%
\begin{tasks}(4)
    \task $m = 1$
    \task $m = -1$
    \task $m = -6$
    \task $m = 6$
\end{tasks}
}

\questionLinesManual{0}{65}{%
Trong không gian $Oxyz$, cho $(P)\colon x + y - 2z + 5 = 0$ và $(Q)\colon 4x + (2 - m)y + mz - 3 = 0$. Tìm tham số $m$ để $(Q)$ vuông góc với $(P)$.
}{%
\begin{tasks}(4)
    \task $m = -3$
    \task $m = -2$
    \task $m = 3$
    \task $m = 2$
\end{tasks}
}

\questionLinesManual{0}{66}{%
Trong không gian $Oxyz$, cho hai mặt phẳng $(\alpha)\colon x + 2y - z - 1 = 0$ và $(\beta)\colon 2x + 4y - mz - 2 = 0$. Tìm $m$ để $(\alpha)$ và $(\beta)$ song song với nhau.
}{%
\begin{tasks}(4)
    \task $m = 1$
    \task Không tồn tại $m$
    \task $m = -2$
    \task $m = 2$
\end{tasks}
}

% =========================================================================
% DẠNG 23: TẬP HỢP ĐIỂM CÁCH ĐỀU / GIAO TUYẾN CHÙM MẶT PHẲNG (CÂU 67 -> 69)
% =========================================================================
\questionLinesManual{0}{67}{%
Trong không gian $Oxyz$, tìm tập hợp các điểm cách đều cặp mặt phẳng song song $4x - y - 2z - 3 = 0$ và $4x - y - 2z - 5 = 0$.
}{%
\begin{tasks}(2)
    \task $4x - y - 2z - 6 = 0$
    \task $4x - y - 2z - 4 = 0$
    \task $4x - y - 2z - 1 = 0$
    \task $4x - y - 2z - 2 = 0$
\end{tasks}
}

\questionLinesManual{0}{68}{%
Trong không gian $Oxyz$, cho hai mặt phẳng $(Q)\colon 3x - y + 4z + 2 = 0$ và $(R)\colon 3x - y + 4z + 8 = 0$. Phương trình mặt phẳng $(P)$ song song và cách đều $(Q)$ và $(R)$ là:
}{%
\begin{tasks}(2)
    \task $3x - y + 4z + 10 = 0$
    \task $3x - y + 4z + 5 = 0$
    \task $3x - y + 4z - 10 = 0$
    \task $3x - y + 4z - 5 = 0$
\end{tasks}
}

\questionLinesManual{0}{69}{%
Trong không gian $Oxyz$, cho hai mặt phẳng $(P)\colon x - 2y - z + 3 = 0$ và $(Q)\colon 2x + y + z - 1 = 0$. Mặt phẳng $(R)$ đi qua điểm $M(1; 1; 1)$ chứa giao tuyến của $(P)$ và $(Q)$ có dạng $m(x - 2y - z + 3) + (2x + y + z - 1) = 0$. Khi đó giá trị của $m$ là:
}{%
\begin{tasks}(4)
    \task $3$
    \task $\dfrac{1}{3}$
    \task $-\dfrac{1}{3}$
    \task $-3$
\end{tasks}
}

% =========================================================================
% DẠNG 24: KHOẢNG CÁCH TỪ ĐIỂM ĐẾN MẶT PHẲNG (CÂU 70 -> 72)
% =========================================================================
\questionLinesManual{0}{70}{%
Trong không gian $Oxyz$, khoảng cách từ điểm $A(1; 0; 0)$ tới mặt phẳng $(P)\colon 2x + 2y - z + 1 = 0$ bằng:
}{%
\begin{tasks}(4)
    \task $3$
    \task $\sqrt{3}$
    \task $9$
    \task $1$
\end{tasks}
}

\questionLinesManual{0}{71}{%
Trong không gian $Oxyz$, khoảng cách từ điểm $A(3; -2; 4)$ đến mặt phẳng tọa độ $(Oxz)$ bằng:
}{%
\begin{tasks}(4)
    \task $4$
    \task $5$
    \task $3$
    \task $2$
\end{tasks}
}

\questionLinesManual{0}{72}{%
Trong không gian $Oxyz$, cho mặt phẳng $(P)\colon 2x - 2y + z - 3 = 0$. Khoảng cách từ điểm $M(-1; 2; 0)$ đến mặt phẳng $(P)$ bằng:
}{%
\begin{tasks}(4)
    \task $-3$
    \task $3$
    \task $2$
    \task $5$
\end{tasks}
}

% =========================================================================
% DẠNG 25: KHOẢNG CÁCH TỪ ĐIỂM CHỨA CĂN THỨC / THAM SỐ (CÂU 73 -> 75)
% =========================================================================
\questionLinesManual{0}{73}{%
Trong không gian $Oxyz$, cho mặt phẳng $(P)\colon 3x + 4y + 2z + 4 = 0$ và điểm $A(1; -2; 3)$. Tính khoảng cách $d$ từ điểm $A$ đến mặt phẳng $(P)$.
}{%
\begin{tasks}(4)
    \task $d = \dfrac{5}{29}$
    \task $d = \dfrac{5}{\sqrt{29}}$
    \task $d = \dfrac{\sqrt{5}}{3}$
    \task $d = \dfrac{5}{9}$
\end{tasks}
}

\questionLinesManual{0}{74}{%
Trong không gian $Oxyz$, cho điểm $M(-1; 2; 1)$ và mặt phẳng $(P)\colon 2x - y + 2z - 4 = 0$. Khoảng cách từ điểm $M$ đến mặt phẳng $(P)$ bằng:
}{%
\begin{tasks}(4)
    \task $\dfrac{2}{3}$
    \task $1$
    \task $2$
    \task $3$
\end{tasks}
}

\questionLinesManual{0}{75}{%
Trong không gian $Oxyz$, cho hai điểm $A(1; 2; 3)$ và $B(3; 4; 4)$. Tìm tất cả các giá trị của tham số $m$ để khoảng cách từ điểm $A$ đến mặt phẳng $2x + y + mz - 1 = 0$ bằng độ dài đoạn thẳng $AB$.
}{%
\begin{tasks}(4)
    \task $m = 2$
    \task $m = -2$
    \task $m = -3$
    \task $m = \pm 2$
\end{tasks}
}

% =========================================================================
% DẠNG 26: TÌM ĐIỂM TRÊN TRỤC CÁCH ĐỀU HAI MẶT PHẲNG (CÂU 76 -> 78)
% =========================================================================
\questionLinesManual{0}{76}{%
Trong không gian $Oxyz$, điểm $M$ thuộc trục $Oy$ và cách đều hai mặt phẳng $(P)\colon x + y - z + 1 = 0$ và $(Q)\colon x - y + z - 5 = 0$ có tọa độ là:
}{%
\begin{tasks}(4)
    \task $M(0; -3; 0)$
    \task $M(0; 3; 0)$
    \task $M(0; -2; 0)$
    \task $M(0; 1; 0)$
\end{tasks}
}

\questionLinesManual{0}{77}{%
Trong không gian $Oxyz$, có bao nhiêu điểm $M$ trên trục tung $Oy$ thỏa mãn $M$ cách đều hai mặt phẳng $(P)\colon x + y - z + 1 = 0$ và $(Q)\colon x - y + z - 5 = 0$?
}{%
\begin{tasks}(4)
    \task $0$
    \task $1$
    \task $2$
    \task $3$
\end{tasks}
}

\questionLinesManual{0}{78}{%
Trong không gian $Oxyz$, tìm trên trục $Oz$ điểm $M$ cách đều điểm $A(2; 3; 4)$ và mặt phẳng $(P)\colon 2x + 3y + z - 17 = 0$.
}{%
\begin{tasks}(4)
    \task $M(0; 0; -3)$
    \task $M(0; 0; 3)$
    \task $M(0; 0; -4)$
    \task $M(0; 0; 4)$
\end{tasks}
}

% =========================================================================
% DẠNG 27: KHOẢNG CÁCH GIỮA HAI MẶT PHẲNG SONG SONG (CÂU 79 -> 81)
% =========================================================================
\questionLinesManual{0}{79}{%
Trong không gian $Oxyz$, cho hai mặt phẳng song song $(P)\colon 2x - y + z = 0$ và $(Q)\colon 2x - y + z - 7 = 0$. Khoảng cách giữa hai mặt phẳng $(P)$ và $(Q)$ bằng:
}{%
\begin{tasks}(4)
    \task $7$
    \task $7\sqrt{6}$
    \task $6\sqrt{7}$
    \task $\dfrac{7}{\sqrt{6}}$
\end{tasks}
}

\questionLinesManual{0}{80}{%
Trong không gian $Oxyz$, khoảng cách giữa hai mặt phẳng $(P)\colon x + 2y + 2z - 8 = 0$ và $(Q)\colon x + 2y + 2z - 4 = 0$ bằng:
}{%
\begin{tasks}(4)
    \task $1$
    \task $\dfrac{4}{3}$
    \task $2$
    \task $\dfrac{7}{3}$
\end{tasks}
}

\questionLinesManual{0}{81}{%
Trong không gian $Oxyz$, khoảng cách giữa hai mặt phẳng $(P)\colon 6x + 3y + 2z - 1 = 0$ và $(Q)\colon x + \dfrac{1}{2}y + \dfrac{1}{3}z + 8 = 0$ bằng:
}{%
\begin{tasks}(4)
    \task $7$
    \task $8$
    \task $9$
    \task $6$
\end{tasks}
}

% =========================================================================
% DẠNG 28: LẬP MP SONG SONG VÀ CÁCH MẶT PHẲNG CHO TRƯỚC 1 KHOẢNG (CÂU 82 -> 84)
% =========================================================================
\questionLinesManual{0}{82}{%
Trong không gian $Oxyz$, cho mặt phẳng $(P)\colon x + 2y + 2z - 10 = 0$. Phương trình các mặt phẳng $(Q)$ song song với $(P)$ và khoảng cách giữa hai mặt phẳng $(P)$ và $(Q)$ bằng $\dfrac{7}{3}$ là:
}{%
\begin{tasks}(2)
    \task $x + 2y + 2z + 3 = 0;\; x + 2y + 2z - 17 = 0$
    \task $x + 2y + 2z - 3 = 0;\; x + 2y + 2z + 17 = 0$
    \task $x + 2y + 2z + 3 = 0;\; x + 2y + 2z + 17 = 0$
    \task $x + 2y + 2z - 3 = 0;\; x + 2y + 2z - 17 = 0$
\end{tasks}
}

\questionLinesManual{0}{83}{%
Trong không gian $Oxyz$, lập phương trình các mặt phẳng song song với mặt phẳng $(\beta)\colon x + y - z + 3 = 0$ và cách $(\beta)$ một khoảng bằng $\sqrt{3}$.
}{%
\begin{tasks}(2)
    \task $x + y - z + 6 = 0;\; x + y - z = 0$
    \task $x + y - z + 6 = 0$
    \task $x - y - z + 6 = 0;\; x - y - z = 0$
    \task $x + y + z + 6 = 0;\; x + y + z = 0$
\end{tasks}
}

\questionLinesManual{0}{84}{%
Trong không gian $Oxyz$, cho mặt phẳng $(P)\colon 2x - 2y + z - 5 = 0$. Viết phương trình mặt phẳng $(Q)$ song song với mặt phẳng $(P)$, cách $(P)$ một khoảng bằng $3$ và cắt trục $Ox$ tại điểm có hoành độ dương.
}{%
\begin{tasks}(2)
    \task $(Q)\colon 2x - 2y + z + 4 = 0$
    \task $(Q)\colon 2x - 2y + z - 14 = 0$
    \task $(Q)\colon 2x - 2y + z - 19 = 0$
    \task $(Q)\colon 2x - 2y + z - 8 = 0$
\end{tasks}
}

% =========================================================================
% DẠNG 29: GÓC GIỮA HAI MẶT PHẲNG CƠ BẢN VÀ THAM SỐ (CÂU 85 -> 87)
% =========================================================================
\questionLinesManual{0}{85}{%
Trong không gian $Oxyz$, cho điểm $H(2; 1; 2)$ là hình chiếu vuông góc của gốc tọa độ $O$ xuống mặt phẳng $(P)$. Số đo góc giữa mặt phẳng $(P)$ và mặt phẳng $(Q)\colon x + y - 11 = 0$ là:
}{%
\begin{tasks}(4)
    \task $60^\circ$
    \task $30^\circ$
    \task $45^\circ$
    \task $90^\circ$
\end{tasks}
}

\questionLinesManual{0}{86}{%
Trong không gian $Oxyz$, biết hình chiếu của gốc tọa độ $O$ lên mặt phẳng $(P)$ là $H(2; -1; -2)$. Số đo góc giữa mặt phẳng $(P)$ với mặt phẳng $(Q)\colon x - y - 5 = 0$ là:
}{%
\begin{tasks}(4)
    \task $30^\circ$
    \task $45^\circ$
    \task $60^\circ$
    \task $90^\circ$
\end{tasks}
}

\questionLinesManual{0}{87}{%
Trong không gian $Oxyz$, cho mặt phẳng $(P)\colon x - 2y + 2z - 5 = 0$ và mặt phẳng $(Q)\colon x + (2m - 1)z + 7 = 0$ với $m$ là tham số. Tìm tất cả các giá trị của $m$ để $(P)$ tạo với $(Q)$ một góc bằng $45^\circ$.
}{%
\begin{tasks}(4)
    \task $\left[\begin{aligned} m &= 1 \\ m &= 4 \end{aligned}\right.$
    \task $\left[\begin{aligned} m &= 2 \\ m &= -2\sqrt{2} \end{aligned}\right.$
    \task $\left[\begin{aligned} m &= 2 \\ m &= 4 \end{aligned}\right.$
    \task $\left[\begin{aligned} m &= 4 \\ m &= \sqrt{2} \end{aligned}\right.$
\end{tasks}
}

% =========================================================================
% DẠNG 30: THỂ TÍCH TỨ DIỆN $OABC$ ĐẠT CỰC TRỊ (CÂU 88 -> 90)
% =========================================================================
\questionLinesManual{0}{88}{%
Trong không gian $Oxyz$, cho mặt phẳng $(P)$ đi qua điểm $M(9; 1; 1)$ cắt các tia $Ox, Oy, Oz$ lần lượt tại $A, B, C$ phân biệt. Thể tích khối tứ diện $OABC$ đạt giá trị nhỏ nhất bằng:
}{%
\begin{tasks}(4)
    \task $\dfrac{81}{2}$
    \task $\dfrac{243}{2}$
    \task $\dfrac{81}{6}$
    \task $243$
\end{tasks}
}

\questionLinesManual{0}{89}{%
Trong không gian $Oxyz$, cho điểm $M(1; 1; 2)$. Mặt phẳng $(P)$ đi qua $M$ cắt các tia $Ox, Oy, Oz$ lần lượt tại $A, B, C$. Thể tích khối tứ diện $OABC$ đạt giá trị nhỏ nhất bằng:
}{%
\begin{tasks}(4)
    \task $\dfrac{9}{2}$
    \task $18$
    \task $9$
    \task $\dfrac{32}{3}$
\end{tasks}
}

\questionLinesManual{0}{90}{%
Trong không gian $Oxyz$, phương trình của mặt phẳng $(P)$ đi qua điểm $M(4; 9; 1)$ và cắt các tia $Ox, Oy, Oz$ lần lượt tại $A, B, C$ sao cho thể tích tứ diện $OABC$ đạt giá trị nhỏ nhất là:
}{%
\begin{tasks}(2)
    \task $9x + 4y + 36z - 108 = 0$
    \task $9x + 4y + 1945z - 2017 = 0$
    \task $-9x + 4y - 36z + 36 = 0$
    \task $9x - 4y + z - 18 = 0$
\end{tasks}
}

% =========================================================================
% DẠNG 31: ĐIỂM CỐ ĐỊNH VÀ HÌNH CHÓP ĐỀU (CÂU 91 -> 93)
% =========================================================================
\questionLinesManual{0}{91}{%
Trong không gian $Oxyz$, cho ba điểm $A(a; 0; 0)$, $B(0; b; 0)$, $C(0; 0; c)$ với $a, b, c$ là ba số thực dương thay đổi thỏa mãn $\dfrac{1}{a} + \dfrac{1}{b} + \dfrac{1}{c} = 2017$. Khi đó mặt phẳng $(ABC)$ luôn đi qua một điểm cố định có tọa độ là:
}{%
\begin{tasks}(2)
    \task $\left(\dfrac{1}{2017}; \dfrac{1}{2017}; \dfrac{1}{2017}\right)$
    \task $(1; 1; 1)$
    \task $\left(\dfrac{1}{3}; \dfrac{1}{3}; \dfrac{1}{3}\right)$
    \task $(2017; 2017; 2017)$
\end{tasks}
}

\questionLinesManual{0}{92}{%
Trong không gian $Oxyz$, cho hai mặt phẳng $(P)\colon x + 4y - 2z - 6 = 0$ và $(Q)\colon x - 2y + 4z - 6 = 0$. Mặt phẳng $(\alpha)$ chứa giao tuyến của $(P), (Q)$ và cắt các trục tọa độ tại $A, B, C$ sao cho hình chóp $O.ABC$ là hình chóp đều. Phương trình của $(\alpha)$ là:
}{%
\begin{tasks}(2)
    \task $x + y + z - 6 = 0$
    \task $x + y + z + 6 = 0$
    \task $x + y + z - 3 = 0$
    \task $x + y - z - 6 = 0$
\end{tasks}
}

\questionLinesManual{0}{93}{%
Trong không gian $Oxyz$, cho các điểm $A(2; 0; 0)$, $B(0; 4; 0)$, $C(0; 0; 6)$ và $D(2; 4; 6)$. Gọi $(P)$ là mặt phẳng song song với $(ABC)$ sao cho $(P)$ cách đều điểm $D$ và mặt phẳng $(ABC)$. Phương trình của $(P)$ là:
}{%
\begin{tasks}(2)
    \task $6x + 3y + 2z - 24 = 0$
    \task $6x + 3y + 2z - 12 = 0$
    \task $6x + 3y + 2z = 0$
    \task $6x + 3y + 2z - 36 = 0$
\end{tasks}
}

% =========================================================================
% DẠNG 32: CỰC TRỊ NGHỊCH ĐẢO ĐỘ DÀI ĐOẠN CHẮN (CÂU 94 -> 96)
% =========================================================================
\questionLinesManual{0}{94}{%
Trong không gian $Oxyz$, viết phương trình mặt phẳng $(P)$ đi qua điểm $M(1; 2; 3)$ và cắt các tia $Ox, Oy, Oz$ lần lượt tại ba điểm $A, B, C$ sao cho biểu thức $\dfrac{1}{OA^2} + \dfrac{1}{OB^2} + \dfrac{1}{OC^2}$ đạt giá trị nhỏ nhất.
}{%
\begin{tasks}(2)
    \task $x + 2y + 3z - 11 = 0$
    \task $x + 2y + 3z - 14 = 0$
    \task $x + 2y + z - 14 = 0$
    \task $x + y + z - 6 = 0$
\end{tasks}
}

\questionLinesManual{0}{95}{%
Trong không gian $Oxyz$, cho ba điểm $A(a; 0; 0)$, $B(0; b; 0)$, $C(0; 0; c)$ với $a, b, c > 0$ thay đổi thỏa mãn $\dfrac{2}{a} - \dfrac{2}{b} + \dfrac{1}{c} = 1$. Khoảng cách từ gốc tọa độ $O$ đến mặt phẳng $(ABC)$ đạt giá trị lớn nhất bằng:
}{%
\begin{tasks}(4)
    \task $3$
    \task $2$
    \task $1$
    \task $4$
\end{tasks}
}

\questionLinesManual{0}{96}{%
Trong không gian $Oxyz$, cho ba điểm $A(a; 0; 0)$, $B(0; b; 0)$, $C(0; 0; c)$ với $a, b, c > 0$. Mặt phẳng $(ABC)$ đi qua điểm $I(1; 2; 3)$ sao cho thể tích tứ diện $OABC$ đạt giá trị nhỏ nhất. Khi đó đẳng thức nào sau đây đúng?
}{%
\begin{tasks}(4)
    \task $a + b + c = 12$
    \task $a^2 + b = c + 6$
    \task $a + b + c = 18$
    \task $a + b - c = 0$
\end{tasks}
}

% =========================================================================
% DẠNG 33: TỐI ƯU HÓA ĐỘ DÀI VÉCTƠ VÀ BẤT ĐẲNG THỨC TAM GIÁC TRÊN MP (CÂU 97 -> 100)
% =========================================================================
\questionLinesManual{0}{97}{%
Trong không gian $Oxyz$, cho ba điểm $A(0; 1; 2)$, $B(1; 1; 1)$, $C(2; -2; 3)$ và mặt phẳng $(P)\colon x - y + z + 3 = 0$. Tìm điểm $M$ thuộc mặt phẳng $(P)$ sao cho $|\vec{MA} + \vec{MB} + \vec{MC}|$ đạt giá trị nhỏ nhất.
}{%
\begin{tasks}(4)
    \task $M(1; 0; 2)$
    \task $M(0; 1; 1)$
    \task $M(-1; 2; 0)$
    \task $M(-3; 1; 1)$
\end{tasks}
}

\questionLinesManual{0}{98}{%
Trong không gian $Oxyz$, cho ba điểm $A(1; 1; 0)$, $B(3; -1; 2)$, $C(-1; 6; 7)$. Tìm tọa độ điểm $M$ thuộc mặt phẳng tọa độ $(Oxz)$ sao cho biểu thức $MA^2 + MB^2 + MC^2$ đạt giá trị nhỏ nhất.
}{%
\begin{tasks}(4)
    \task $M(3; 0; -1)$
    \task $M(1; 0; 0)$
    \task $M(1; 0; 3)$
    \task $M(1; 1; 3)$
\end{tasks}
}

\questionLinesManual{0}{99}{%
Trong không gian $Oxyz$, cho hai điểm $A(1; -1; 1)$, $B(0; 1; -2)$ và điểm $M$ di động trên mặt phẳng tọa độ $(Oxy)$. Giá trị lớn nhất của biểu thức $T = |MA - MB|$ bằng:
}{%
\begin{tasks}(4)
    \task $\sqrt{6}$
    \task $\sqrt{12}$
    \task $\sqrt{14}$
    \task $\sqrt{8}$
\end{tasks}
}

\questionLinesManual{0}{100}{%
Trong không gian $Oxyz$, cho bốn điểm $A(1; -2; 0)$, $B(0; -1; 1)$, $C(2; 1; -1)$, $D(3; 1; 4)$. Hỏi có tất cả bao nhiêu mặt phẳng cách đều cả bốn điểm đã cho?
}{%
\begin{tasks}(4)
    \task $1$
    \task $4$
    \task $7$
    \task Vô số
\end{tasks}
}

\end{document}